\documentclass[11pt, a4paper, twoside]{article}
\usepackage[T1]{fontenc}
\usepackage[utf8]{inputenc}
\usepackage{amssymb,amsmath}
\usepackage[portuguese]{babel}
\usepackage{comment}
\usepackage{datetime}
\usepackage[pdfusetitle]{hyperref}
\usepackage[all]{xy}
\usepackage{graphicx}
\addtolength{\parskip}{.5\baselineskip}

%aqui comeca o que eu fiz de verdade, o resto veio e eu to com medo de tirar
\usepackage{xcolor}
\usepackage{listings} %biblioteca pro codigo
\usepackage{color}    %deixa o codigo colorido bonitinho
\usepackage[landscape, left=0.5cm, right=0.5cm, top=1cm, bottom=1.5cm]{geometry} %pra deixar a margem do jeito que o brasil gosta

\definecolor{gray}{rgb}{0.4, 0.4, 0.4} %cor pros comentarios
%\renewcommand{\footnotesize}{\small} %isso eh pra mudar o tamanho da fonte do codigo
\setlength{\columnseprule}{0.2pt} %barra separando as duas colunas
\setlength{\columnsep}{10pt} %distancia do texto ate a barra

\lstset{ %opcoes pro codigo
breaklines=true,
keywordstyle=\color{blue}\bfseries,
commentstyle=\color{gray},
breakatwhitespace=true,
language=C++,
%frame=single, % nao sei se gosto disso ou nao
numbers=none,
rulecolor=\color{black},
showstringspaces=false
stringstyle=\color{blue},
tabsize=4,
basicstyle=\ttfamily\footnotesize, % fonte
}
\lstset{literate=
%   *{0}{{{\color{red!20!violet}0}}}1
%    {1}{{{\color{red!20!violet}1}}}1
%    {2}{{{\color{red!20!violet}2}}}1
%    {3}{{{\color{red!20!violet}3}}}1
%    {4}{{{\color{red!20!violet}4}}}1
%    {5}{{{\color{red!20!violet}5}}}1
%    {6}{{{\color{red!20!violet}6}}}1
%    {7}{{{\color{red!20!violet}7}}}1
%    {8}{{{\color{red!20!violet}8}}}1
%    {9}{{{\color{red!20!violet}9}}}1
%	 {l}{$\text{l}$}1
	{~}{$\sim$}{1} % ~ bonitinho
}

\title{身正不怕影子斜 \\ Unioeste}
\author{Amós, Carlos e Eduarda}
\phrase{Um pé reto não tem medo de um sapato torto.}


\begin{document}
\twocolumn
\date{\today}
\maketitle


\renewcommand{\contentsname}{Índice} %troca o nome do indice para indice
\tableofcontents


%%%%%%%%%%%%%%%%%%%%
%
% DP
%
%%%%%%%%%%%%%%%%%%%%

\section{DP}

\subsection{Bitmask DP}
\begin{lstlisting}
//hamiltonian path
9a4 const int MAXN = 20;
2bd const int MAXM = (1<<20)+7;
3e4 const int MOD = 1e9+7;
47b vector<int> graph[MAXN];
50a int dp[MAXN][MAXM];
1a8 int n; 
 
003 int pd(int v, int mask){
718 	auto& p = dp[v][mask];
    	//cout << v+1 << " " << mask << endl;
6d2 	if(v == n-1){
    		//cout << v+1 << " " << mask << endl;
2bf 		return p = (mask ==((1<<n)-1));
c4a 	}
     
631 	if(p != -1) return p;
     
1a4 	int ans = 0;
    	//cout << graph[v].size() << endl;
284 	for(auto u : graph[v]){
    		//cout << u+1 << " " << mask << " " << (1<<u) << endl;
f67 		if(!((mask>>u)&1)){
c1c 			int newmask = mask | (1 << u);
0bc 			ans = (ans + pd(u, newmask))%MOD;
58a 		}
d75 	}
46a 	return p = ans;
b5a }

//iterativo
199 dp[1][0] = 1;
51c for (int s = 2; s < 1 << city_num; s++) {
    	// only consider subsets that have the first city
ba0 	if ((s & (1 << 0)) == 0) continue;
    	// also only consider subsets with the last city if it's the full subset
8fa 	if ((s & (1 << (city_num - 1))) && s != ((1 << city_num) - 1)) continue;
631 	for (int end = 0; end < city_num; end++) {
ec6 		if ((s & (1 << end)) == 0) continue;
    		// the subset that doesn't include the current end
171 		int prev = s - (1 << end);
7de 		for (int j : come_from[end]) {
fc6 			if ((s & (1 << j))) {
597 				dp[s][end] += dp[prev][j];
9f8 				dp[s][end] %= MOD;
7c9 			}
9be 		}
9ff 	}
290 }

//merging subsets O(n^3)
827 for (int mask = 0; mask < (1 << n); mask++) {
e0a 	for (int submask = mask; submask != 0; submask = (submask - 1) & mask) {
982 		int subset = mask ^ submask;
    		// do whatever you need to do here
70a 	}
e71 }

//submasks
199 dp[1][0] = 1;
51c for (int s = 2; s < 1 << city_num; s++) {
    	// only consider subsets that have the first city
ba0 	if ((s & (1 << 0)) == 0) continue;
    	// also only consider subsets with the last city if it's the full subset
8fa 	if ((s & (1 << (city_num - 1))) && s != ((1 << city_num) - 1)) continue;
631 	for (int end = 0; end < city_num; end++) {
ec6 		if ((s & (1 << end)) == 0) continue;
    		// the subset that doesn't include the current end
171 		int prev = s - (1 << end);
7de 		for (int j : come_from[end]) {
fc6 			if ((s & (1 << j))) {
597 				dp[s][end] += dp[prev][j];
9f8 				dp[s][end] %= MOD;
7c9 			}
9be 		}
9ff 	}
290 }
\end{lstlisting}

\subsection{Digit DP}
\begin{lstlisting}
// TEMPLATE 1: RECURSIVA Padrão [0, R] ou [1, R]
// ==========================================
730 int dp[20][11][2][2]; // [index][state][tight][ldz]

c3f int pd(int index, int state, int tight, int ldz) {
2e8     if(index == n) {
505         return /* condicao base */;
d0c     }
        
890     if (dp[index][state][tight][ldz] != -1) return dp[index][state][tight][ldz];
    
c23     int ub = tight ? r[index] - '0' : 9;
1a4     int ans = 0;
    
70b     for (int digit = 0; digit <= ub; digit++) {
            // if(!ldz && condicao_invalida) continue;
    
e16         int new_tight = tight && (digit == ub);
d3b         int new_ldz = ldz && (digit == 0);
            
f1b         int new_state = state; // atualiza seu estado
    
cf0         ans += pd(index + 1, new_state, new_tight, new_ldz);
572     }
c77     return dp[index][state][tight][ldz] = ans; 
18b }

// ==========================================
// TEMPLATE 2: RECURSIVA Range [L, R] Simultâneo
// ==========================================
377 int dp[20][11][2][2][2]; // [index][state][above][under][ldz]

d6c int pd(int index, int state, int above, int under, int ldz) {
